\subsection*{How often to poll}

Our first plan was to poll every 15 minutes. Before we started, we tested how long a finished
trip stays in the feed: on 26 October we saved 32 snapshots, 5 minutes apart, and compared the
last snapshot of each trip with the time of its last stop (\code{scripts/retention\_check.py}).
Of 771 trips that left the feed, all stayed between 4 and 11 minutes after the last stop
(median 8), and none stayed 15. So 15 minutes would have been too slow: for most trips we would
never see the real time at the last stops.

With 5 minutes, almost every trip is seen after it has finished. A trip that leaves after only
4 minutes can still fall between two polls, and then its last stops keep the predicted time
instead of the real one. We count these stops in every Silver run: 0.2--1.8\% a day. Leaving
them out changes the on-time share by only 0.05 points, so we keep them. Polling more often
would save these few stops but multiply storage and API calls.

\subsection*{How the raw data is kept}

\begin{itemize}
  \item Every answer is saved exactly as received, gzipped (about 15 times smaller, and Spark
        reads it directly), in folders by UTC date and hour.
  \item Every run, also a failed one, adds one line to a metadata file: time, HTTP status, size,
        checksum and error. Broken XML is saved too and marked as broken.
  \item A file is never overwritten, and every Silver row keeps the id of its Bronze file, so
        any number in Gold can be traced back to one raw answer.
\end{itemize}

Bronze runs from cron on the prod VM, not from a GitHub schedule. When we tried both, every
snapshot was fetched twice, and GitHub's scheduled runs started minutes late, which is enough
for a trip to leave the feed. From 27 October to 22 November there were 7{,}778 runs, one
failed (a DNS error on 29 October), and there was no gap longer than 10 minutes.

\subsection*{Line names come from Bronze too}

The feed only has line ids like \code{KOL:Line:8\_2043}. Entur's Journey Planner~\cite{jp}
knows the public number (here X39). Our first version called this API and wrote the answer
directly into a table, so the raw answer was never kept. We got this feedback at the
presentation and changed it:
\begin{itemize}
  \item The step \code{bronze-lines} saves the full Kolumbus line list (320 lines) once a day in
        its own Bronze folder, raw, with a metadata line, the same way as the SIRI data.
  \item The step \code{lines} reads all saved lists with Spark and keeps the newest name for
        each line id. A line that stops running keeps its last name.
  \item If the API is down, we only get a warning and Gold uses the lists already saved.
\end{itemize}

\begin{lstlisting}
w = Window.partitionBy("line") \
          .orderBy(F.col("fetched_at").desc())
names = lines.withColumn("_rn", F.row_number().over(w)) \
             .filter("_rn = 1").drop("_rn")
\end{lstlisting}
